\section{Floating Point Unit (FPU)}
\begin{itemize}
    \item Program Control \& Status Register (FPCSR)
    \item Grundoperationen: Mathematische Befehle, Datentransferbefehle, Konvertierungsbefehle, Compare-Befehl
\end{itemize}

\smallbreak
\textbf{Register FPU}

\begin{itemize}
    \item 32 Single Precision Register (S0-S31)
    \item 16 Double Precision Shadow Register (D0-D15)
\end{itemize}

\subsection{IEEE 754: Addition / Subtraktion}
Nötige Schritte:
\begin{enumerate}
    \item Beide Werte müssen normalisiert sein.
    \item Exponenten auf grösseren Wert anpassen.
    \item Mantissen addieren bzw. subtrahieren.
    \item Resultat wenn nötig normalisieren.
    \item Umwandlung in IEEE 754 Format.
\end{enumerate}
Beispiel: $0.74 + 0.09 = 0.83$ \\
\textit{Runden:} Ergebnis auf nächstliegende Zahl runden (round to nearest). \\
\textit{Problematik:} Bei extremer Anpassung der Exponenten wird der Rundungsfehler immer massiver.

\vspace{0.5em}
\noindent\framebox{\parbox{\linewidth}{\centering \vspace{10pt}\textit{[Placeholder: Ausführliche Binär-Berechnung (Addition $0.74+0.09$)]}\vspace{10pt}}}
\vspace{0.5em}

\subsection{IEEE 754: Multiplikation / Division}
Nötige Schritte:
\begin{enumerate}
    \item Überprüfung auf Werte mit 0.0f
    \item Exponenten addieren bzw. subtrahieren
    \item Mantissen multiplizieren bzw. dividieren und Vorzeichen anpassen
    \item Resultat wenn nötig normalisieren
    \item Umwandlung in IEEE 754 Format
\end{enumerate}
Beispiel: $6.625 \cdot (-0.5) = -3.3125$ \\
\textit{Runden:} Ebenfalls "round to nearest".

\vspace{0.5em}
\noindent\framebox{\parbox{\linewidth}{\centering \vspace{10pt}\textit{[Placeholder: Ausführliche Binär-Berechnung (Multiplikation $6.625 \cdot -0.5$)]}\vspace{10pt}}}
\vspace{0.5em}